\section{Discussion}
\label{sec:discussion}

\subsection{Do OBTF Load and Renewable Forecasts Improve Price Forecasts?}
\label{sec:discussion_rq2}

The OBTF forecasts improve the price forecast. Replacing the
published \mbox{ENTSO-E} load forecast with the OBTF load, combined
renewable-generation, and residual-load forecasts lowers the day-ahead price
RMSE by 11.2\% relative to the base price forecast, and the improvement is
statistically significant. Both forecasts use the same estimator, the same
forecast creation time, and the same price-lag, calendar, and raw-weather
inputs. They do not, however, share the same information set. The OBTF models
draw on a much richer weather representation, installed capacity from the
MaStR, population weights, and realized load, which the base
model does not receive directly, because hundreds of such features could not be
estimated from 70 daily training observations. The gain therefore measures the
value of this information delivered in the compact form of explicit forecasts,
not the value of the forecast representation alone
(Section~\ref{sec:limitations}).
The OBTF load and renewable-generation forecasts thus carry
predictive information that the price model does not recover on its own from
price history, raw weather, and the public load forecast.
This agrees with \textcite{maciejowska2021enhancing}, who show that including
load, wind, and solar forecasts, and improving their accuracy, yields more
accurate day-ahead price forecasts. Their comparison is against price models
without these forecasts, whereas the base forecast here already contains raw
weather, so the result extends their finding to a price model that already has
access to the underlying weather information.

The weather-free variant refines this interpretation. A price model that keeps
the OBTF forecasts but drops the raw weather features is statistically
indistinguishable from the model that retains both, and it still improves on the
base forecast by essentially the same margin. Once the OBTF forecasts
are included, the two-cluster raw weather features thus add no detectable
accuracy. This is consistent with how the forecasts are constructed, because
each is produced by a task-specific model that already distills the relevant
weather signal into an expected load, solar, or wind quantity. For the
data-limited price model, these few summary variables appear to be a more usable
form of the weather information than the raw weather from which they are
derived, although a non-significant difference does not prove that the raw
weather is redundant.

The most operationally relevant finding emerges when this result is read together
with RQ1. The OBTF solar and wind forecasts are less accurate than the
corresponding \mbox{ENTSO-E} benchmarks, yet they still improve the price
forecast, whereas the more accurate \mbox{ENTSO-E} renewable-generation forecasts,
which several price studies use as direct inputs
\cite{trebbien2023probabilistic,marcjasz2023distributional,hilger2024multivariate},
cannot be used at all because they are published only after gate closure. For
an operational price forecast, accuracy and admissibility are therefore both
required. A more accurate benchmark that is not observable before the deadline
cannot serve as an input, and an admissible forecast is useful only if it adds
information beyond price history, weather, and the public load forecast. The
OBTF forecasts satisfy both conditions, and the OBTF load forecast additionally exceeds the accuracy of the public forecast it is compared to.

\subsection{Operational Value of Waiting}

RQ3 shows that the accuracy of the operational price forecast improves as the
morning progresses and the admissible information set grows, but gradually and,
before 12:00, without statistically significant steps. Between 07:00 and 09:00
the forecast barely changes. From 10:00, when the aFRR results become available, and 11:00, when the public load
forecast and the mFRR results enter, the error
declines more clearly and is 5.6\% lower at 11:00 than at 07:00. On the
six-month test period, however, these gains are not statistically significant,
which may partly reflect that the backtest uses the same weather run at all
cutoffs (Section~\ref{sec:limitations}). A
participant who must commit early therefore loses little by forecasting at 07:00
rather than waiting until 11:00, although the direction of the evidence favors
waiting.

A further and larger gain comes at the gate-closure cutoff, once the EXAA price
is admissible. The EXAA auction clears earlier than the SDAC auction and its
price is a direct market estimate of the same day-ahead quantity, so it is
unsurprising that it carries more short-horizon information than any fundamental
input, consistent with the predictive value of the EXAA auction documented by
\textcite{ziel2015exaa}. That a compact model using the EXAA price alone matches
the full specification is a known property of this market rather than a finding
of this thesis \cite{eiser2026operational}, and the feature-importance analysis
reflects it, with the EXAA price accounting for about two-thirds of the
contribution mass at 12:00.

This does not diminish the value of the OBTF forecasts but clarifies
where each source matters. The EXAA price is published only shortly before gate
closure, so throughout the earlier part of the morning, while a participant is
still forming a bid rather than submitting it at the deadline, no market price is
yet available. Across those cutoffs the forecast rests on price lags, weather, the OBTF load and renewable-generation forecasts, and the reserve-market results,
which lower the error at every cutoff before 12:00, although not by a robustly
significant margin. Once the EXAA price arrives, it dominates.

\subsection{Role of Weather and Public Inputs}

The literature review left open how weather is best represented in a day-ahead
price model. The weather-free variant and the feature-importance results help
answer this question. In the base model,
raw weather is one of the two largest groups of inputs, so the fitted model
relies heavily on weather. When the OBTF
forecasts are added, the importance of raw weather drops sharply and is taken over by the OBTF residual-load and renewable-generation forecasts, and the
weather-free variant shows that the raw weather features can be removed
altogether without a measurable loss. For this data-limited price model, the
OBTF forecasts are therefore at least as useful a representation of
the weather-driven system state as the two-cluster raw weather fields, plausibly
because they condense a much richer weather signal into a small number of
variables that are already aligned with how weather affects prices. The dominance of the OBTF residual-load forecast is consistent with residual demand as a
central driver of the price \cite{do2021residualDemand,pape2016fundamentals} and
with the explainable machine-learning analysis of \textcite{trebbien2023merit},
who find load, wind, and solar generation to be the most important features for
German day-ahead prices. This
contrasts with operational pipelines that instead feed spatially resolved
weather directly into the price model in place of the unavailable renewable-generation forecasts \cite{eiser2026operational}, and with studies that select among a large
set of weather variables within the price model \cite{ludwig2015bigdata}. Meteorological
forecasts have been shown to improve German price forecasts beyond the day-ahead
horizon \cite{sgarlato2023weatherPredictions}. At the day-ahead horizon, the
results here suggest that explicitly
forecasting load and renewable generation is a promising way to bring weather
information into a day-ahead price model, subject to the caveat on the unequal
weather information discussed in the limitations.

The public data sources play distinct and complementary roles. The published
\mbox{ENTSO-E} load forecast is a useful baseline whose accuracy the
OBTF load model improves on by correcting it with later information, so
in this setting the public forecast is more useful as a starting point for
correction than as a final input. The reserve-market results add some information about system tightness
before any market price is available, although their effect is not robust on
this test period. The EXAA price is the single most informative
public input at gate closure, but it is admissible only at the very end of the
morning. Read together, these inputs form a sequence in which fundamental and
reserve information dominate early and the exchange price dominates late, and an
operational price model benefits from using each of them at the point in the
morning where it becomes available.

\subsection{Openness and Operational Deployment}

A defining property of the pipeline is that it is fully open and operationally
admissible. Every input is drawn from open, automatically retrievable sources,
namely open weather forecasts, the \mbox{ENTSO-E} Transparency Platform, the
MaStR, and gridded population data, so the self-generated load,
solar, wind, and price forecasts can be generated and extended without
proprietary inputs. The preprocessing, the OBTF models, and the price model
are released as a modular, configuration-driven codebase that is ready for
operational use and deployment, so each stage can be reused and extended.
Together with the daily Energy-Arena submission, this follows the call of
\textcite{lago2021forecasting} for open-access benchmarks in electricity price
forecasting. By replacing the
post-gate-closure \mbox{ENTSO-E} renewable-generation forecasts with
OBTF forecasts, the pipeline also respects the operational timing,
so the reported accuracy reflects what a participant could actually have
produced before the deadline.

The pipeline is additionally submitted to Energy-Arena \cite{kleinebrahm2026energyarena}, where its load, solar, wind,\footnote{The Energy-Arena
\mbox{DE-LU} wind challenge targets onshore wind generation, so only the onshore
component of the OBTF wind forecast is submitted there, whereas the OBTF wind forecast evaluated in this thesis is total wind generation, onshore and
offshore.} and price forecasts are
scored against realized outcomes under a common gate-closure deadline and shared
in the benchmark's cooperative mode so that other forecasters can reuse them, for
example as ensemble members. This provides a continuous, forward-looking, and
leakage-controlled evaluation that extends the fixed test period of this thesis
and lets the OBTF forecasts serve as a reusable building block for further
work.

Beyond serving the price model, the OBTF component forecasts are
directly reusable. Regenerated from the open repository, the OBTF load forecast improves on the public \mbox{ENTSO-E} load forecast and can be used in its place,
while the OBTF solar and wind forecasts provide admissible substitutes for the
\mbox{ENTSO-E} renewable-generation forecasts, which are not published before
gate closure. Unlike the \mbox{ENTSO-E} forecasts, which are produced by
undisclosed models and cannot be inspected, the OBTF forecasts are
fully transparent, because their inputs, models, and construction are open and
can be examined and adapted rather than taken as given.

\subsection{Limitations}
\label{sec:limitations}

The load benchmarking in RQ1 is an operational rather than a like-for-like
modeling comparison. The OBTF load forecast corrects the public
\mbox{ENTSO-E} forecast with information available up to the pre-gate-closure
cutoff, whereas the \mbox{ENTSO-E} forecast is produced by an undisclosed model
at a generally earlier, unspecified cutoff. A comparison that fixes both
forecasts to an identical cutoff is therefore not feasible, and the reported load
improvement should be read as the operational value of correcting the public
forecast before gate closure rather than as evidence of a superior standalone
load model. This caveat concerns only the interpretation of the standalone load
comparison. Because the OBTF load forecast serves solely as an admissible input to the
price model, its only operational requirement is to be available before the 12:00
gate closure, so producing it at a later cutoff than the \mbox{ENTSO-E} forecast,
with correspondingly more recent information, is by design rather than a
disadvantage.

The central caveat of the RQ2 comparison, discussed in
Section~\ref{sec:discussion_rq2}, is that the base price model and the OBTF models do
not share the same information set. The base model receives raw weather only
through a compact two-cluster aggregation of the \mbox{ICON-D2} fields and none
of the capacity, population, and realized-load inputs of the OBTF models. A
comparison that supplied the base model with all of these inputs directly would
isolate the effect of the explicit representation, but hundreds of such features
cannot be estimated reliably from 70 daily training observations. The OBTF forecasts avoid this by splitting the
problem. Load, solar, and wind generation depend on different weather
variables, such as temperature for load, irradiance and cloud cover for solar,
and wind speed for wind. Without the first-stage forecasts, a single price model
would have to receive all of these variables at once. With them, each OBTF model uses only the weather it needs, and the price model receives the result as
a small number of forecast series, namely load, renewable generation, and
residual load. This limitation is thus at the same time the central appeal of
the approach.

The evaluation covers a single test window from February to July 2026, which is
the period for which post-transition quarter-hourly prices and the full set of
operational inputs are available. The results therefore describe one half-year of
a market that only recently moved to quarter-hourly clearing, and the autumn and
mid-winter months lie outside the test window. Whether the reported effects are
stable across seasons and across years with different fuel-price and
renewable-capacity conditions cannot be established from this sample, and a longer
multi-year evaluation would be needed to settle it. The ongoing Energy-Arena submission
extends this evidence with each new delivery day.

A major limitation of RQ3 concerns the weather input. The earlier
\mbox{ICON-D2} runs that would be admissible at the 07:00 to 09:00 cutoffs are
not archived for the whole test period, so all cutoffs use the 06~UTC run,
which is available only from about 09:20. The early cutoffs therefore receive a
more recent weather forecast than they could use in operation, which favors
them. At the same time, the backtest cannot capture any improvement from newer
weather runs over the morning, although weather is a central input of both the
OBTF and the price models. Both effects bias the measured gain from waiting
downward and may explain why the gains before 12:00 are not statistically
significant. A backtest with the operational run schedule of
Table~\ref{tab:cutoff_accrual}, possible once the earlier runs have been archived
over a full test period, would resolve this.

The OBTF solar and wind forecasts remain less accurate than the
\mbox{ENTSO-E} benchmarks, and the feature-importance analysis attributes more of
the price model's contribution to the OBTF load and residual-load forecasts than to the OBTF renewable-generation forecast (Section~\ref{sec:results_feature_importance}). The
OBTF solar and wind forecasts are thus admissible and useful but not competitive
in absolute accuracy, and closing this gap is a limitation of the current OBTF models rather than a property of the price model. The capacity data from the
current MaStR also include late-registered units that were not
yet known at the time, an effect the rolling bias correction partly absorbs. In
addition, the archived \mbox{ICON-D2} irradiance is incomplete or implausible on a
few days in June 2026. These days are retained, so the affected solar forecasts,
and through them the price forecasts, carry these input errors. As with the OBTF load forecast, however, the comparison is not on equal footing. The \mbox{ENTSO-E}
solar and wind benchmarks are produced by undisclosed models and, being published
only after gate closure, may draw on later weather runs and more recent
information than the OBTF forecasts, which are restricted to the
pre-gate-closure cutoff. Part of the observed gap may therefore reflect this
information and timing advantage rather than modeling skill alone, and how much of
it is genuinely attributable to the OBTF models cannot be established from
this comparison.

The feature-importance analysis is diagnostic rather than causal. The SHAP
contributions describe how the fitted model distributes its predictions across
inputs on the sampled data, not the true causal effect of each input, and they
are computed on the variance-stabilized model target so that only their direction
and relative magnitude are interpretable.

Finally, the thesis evaluates point forecasts only. Operational trading and risk
management typically require the full predictive distribution, and the pipeline
does not yet produce or evaluate probabilistic price forecasts.

\subsection{Future Work}

Several extensions follow directly from these limitations. The most immediate is
to close the accuracy gap of the OBTF solar and wind forecasts, for
example through multi-model weather ensembles \cite{bruninx2026windensembles}, wind speeds interpolated to the
turbine-specific hub heights recorded in the MaStR, or an
explicit treatment of curtailment, since more accurate renewable-generation forecasts would
strengthen the fundamental contribution to the price model.

The estimators themselves could also be revisited. Tabular
foundation models such as TabPFN are designed for small training sets
\cite{hollmann2025tabpfn}, which makes them a natural candidate for the
data-limited price model with its 70-day training window.

Beyond point forecasts, the pipeline could be extended to probabilistic
forecasting. The OBTF models could provide quantile forecasts of load, solar, and
wind generation as price-model inputs, which \textcite{uniejewski2026probabilistic}
show to lower the price RMSE by up to 13\% compared with point inputs for the
German market. The same OBTF forecasts and cutoff structure could feed a
distributional or quantile price model \cite{marcjasz2023distributional}, which would make the operational value of
the approach measurable in terms of the full predictive distribution rather than
the conditional mean alone.

Further gains may come from a broader admissible information set. Cross-border
prices and residual-load forecasts of neighboring bidding zones, as well as fuel
and emission-allowance prices, are known to carry predictive information for
German day-ahead prices
\cite{trebbien2023probabilistic,hilger2024multivariate,pape2016fundamentals},
and incorporating them within the open-data and
pre-gate-closure constraints of this thesis would show how much further the
operational forecast can be improved.

Finally, a longer and multi-seasonal evaluation would establish the stability of
the reported effects across market regimes, and extending the setup to the
intraday markets that follow the day-ahead auction would broaden its operational
relevance.
